\chapter{Analyse et conception}
\label{chap:analyse}

\section{Introduction}
\label{sec:analyse:intro}

\section{Étude de l'existant}
\label{sec:analyse:existant}

% Avantages et limites des solutions existantes

\section{Étude des besoins fonctionnels}
\label{sec:analyse:fonctionnels}

\section{Étude des besoins non fonctionnels}
\label{sec:analyse:non-fonctionnels}

% Performance, sécurité, disponibilité, ergonomie, maintenabilité

\section{Identification des acteurs}
\label{sec:analyse:acteurs}

\section{Modélisation du système}
\label{sec:analyse:modelisation}

\subsection{Diagramme de cas d'utilisation}
\label{subsec:analyse:uc}

\begin{figure}[H]
\centering
% \includegraphics[width=\textwidth]{images/diagramme-cas-utilisation.png}
\caption{Diagramme de cas d'utilisation}
\label{fig:uc}
\end{figure}

\subsection{Diagramme de classes}
\label{subsec:analyse:classes}

\begin{figure}[H]
\centering
% \includegraphics[width=\textwidth]{images/diagramme-classes.png}
\caption{Diagramme de classes}
\label{fig:classes}
\end{figure}

\subsection{Diagramme de séquence}
\label{subsec:analyse:sequence}

% Optionnel

\begin{figure}[H]
\centering
% \includegraphics[width=\textwidth]{images/diagramme-sequence.png}
\caption{Diagramme de séquence}
\label{fig:sequence}
\end{figure}

\subsection{Modèle relationnel de la base de données}
\label{subsec:analyse:bdd}

\begin{figure}[H]
\centering
% \includegraphics[width=\textwidth]{images/modele-relationnel.png}
\caption{Modèle relationnel de la base de données}
\label{fig:modele-relationnel}
\end{figure}

\section{Conclusion}
\label{sec:analyse:conclusion}
